%% preface.tex — JA4proxy Operator's User Guide
%% Preface

\chapter*{Preface}
\addcontentsline{toc}{chapter}{Preface}

\newthought{This guide is for the engineer who has to make JA4proxy work.}
Not for the evaluator deciding whether to buy it, and not for the developer
reading the source code. For the security engineer or operations team who has
been handed a repository URL and told to get it in front of production traffic.

If that is you, this document covers what you need: installing the stack, understanding
what the proxy is doing and why, configuring it for your environment, keeping it running,
and responding when something goes wrong.

\section*{How This Guide Is Organised}

The chapters follow the natural sequence of operating \ja\ for the first time and
then day-to-day:

\begin{description}[leftmargin=3.5cm, style=nextline]
  \item[Chapter 1 — Introduction] What \ja\ actually does, and the key concepts
    (TLS fingerprinting, the dial, the pipeline) you need to understand before
    touching any configuration.
  \item[Chapter 2 — Installation] Getting the stack running from scratch. Prerequisites,
    step-by-step setup, and how to verify everything is healthy.
  \item[Chapter 3 — First Steps] What to look at once the proxy is running. Generating
    test traffic, reading the logs, and opening the Grafana dashboard for the first time.
  \item[Chapter 4 — Configuration] The full configuration file walkthrough: the dial,
    country filtering, fingerprint lists, rate limiting, and bypass conditions.
  \item[Chapter 5 — Day-to-Day Operations] Starting and stopping, adding entries to
    lists, managing bans, updating the GeoIP database.
  \item[Chapter 6 — Monitoring] The Grafana dashboard in detail, key Prometheus
    metrics, Loki log queries, and pre-configured alerts.
  \item[Chapter 7 — Incident Response] Detecting an active attack, emergency blocking
    procedures, investigating suspicious fingerprints, recovering from false positives.
  \item[Chapter 8 — Scaling] Adding proxy instances, what state is shared, Kubernetes
    deployment and autoscaling.
  \item[Chapter 9 — Troubleshooting] Common failure modes and how to fix them.
    Diagnostic commands and log analysis.
  \item[Quick Reference] All \code{make} targets, Redis CLI commands, port numbers,
    and metric names on one page.
\end{description}

\section*{The Technical Reference Manual}

This guide deliberately stops short of exhaustive technical detail. Signal definitions,
full metric listings, Redis key schemas, API specifications, and architecture decision
records live in the companion \emph{Technical Reference Manual}. Where this guide says
\emph{see the Reference Manual}, that is where to look.

\section*{Status Note}

\ja\ is a proof-of-concept: functional, tested against real traffic profiles, and
used in pre-production environments. It is \emph{not} yet production-hardened in the
sense that an audited commercial product would be. Secrets management, TLS mutual
authentication between internal services, and several other hardening items are
documented as known gaps in the Reference Manual's \emph{DMZ Deployment Readiness}
chapter. Read that chapter before placing \ja\ in front of live user traffic.

This guide is honest about those limitations. Where something requires additional
hardening before production, it says so plainly.

\section*{Conventions}

\begin{description}[leftmargin=3cm, style=nextline]
  \item[\code{monospace}] Commands, file paths, configuration keys, and code.
  \item[\configkey{config.key}] Configuration file keys in \code{config/proxy.yml}.
  \item[\emph{Side notes}] Supplementary information in the margin.\sidenote{Margin
    notes like this one provide context or cross-references without interrupting
    the main text.}
  \item[\emph{Note boxes}] Provide important information you should read.
  \item[\emph{Warning boxes}] Flag things that can cause operational problems.
  \item[\emph{Danger boxes}] Describe actions that can cause data loss or outages.
\end{description}

\vspace{0.5\baselineskip}
\begin{notebox}[Example note box]
This is a note box. Use these for important context that does not interrupt the main flow.
\end{notebox}

\begin{warningbox}[Example warning box]
This is a warning box. Use these for operational risks.
\end{warningbox}

\vspace{\baselineskip}
\noindent\emph{Seán Ó Ríordáin\\
JA4proxy Project}
